\documentstyle [11pt]{article}
\title{Radical Decomposition and Geometry Theorem Proving}
\author{Michael Kalkbrener\\Department of Mathematics\\
Swiss Federal Institute of 
Technology\\
CH-8092 Zurich, Switzerland}
\date{} 
\begin{document}
\pagestyle{empty}
\maketitle

\vskip 1cm
\centerline{\bf Abstract}

\medskip
Often geometric statements are true only in a generic sense, i.e.
after certain degenerate situations have been ruled out. The problem of
finding nondegeneracy conditions is related to the following algorithmic
problem: for a radical $J$ in a Noetherian commutative ring $R$ with identity
and a polynomial $f$ in $R$ compute radicals $J_1,J_2$ such that 
\[ J=J_1\cap J_2,\ \ \ \ \ \ f\in J_1,\ \ \ \ \ \  J_2:f=J_2. \]
We show that
if this ``splitting'' problem can be solved in $R$ then it can be solved in
$R[x_1,\ldots,x_n]$. From the constructive proof
we obtain an algorithm for splitting radicals in multivariate polynomial
rings over fields.
The basic operation in this algorithm is the 
computation of greatest common divisors of univariate polynomials over 
extension fields. 
No factorization or Gr\"obner bases computation is
needed.
We apply this algorithm to geometry theorem proving and show
that our algorithm can not only be used for deciding correctness
of geometric
statements but
also for constructing 
simplest nondegeneracy conditions with respect
to lexicographical degree orderings.
\end{document}
